Selenium nanoparticles (SeNPs) have attracted widespread attention due to their favorable biocompatibility and diverse bioactivities; however, their cytotoxic responses are co-regulated by particle size, concentration, surface functionalization, and material chemical states, rendering systematic evaluation via conventional experimental assays labor-intensive and low-throughput. To address the challenge of predicting selenium nanoparticle cytotoxicity, this study establishes a two-tiered research framework: "benchmark dataset model comparison + dedicated dataset enhanced-feature modeling." First, leveraging a publicly available multi-type nanoparticle dataset curated by Masarkar et al. \cite{masarkar2024}, a systematic performance comparison was conducted across Random Forest, XGBoost, Attention-Based Random Forest (ABRF), Leaf Attention-Based Random Forest (LARF), TabTransformer, and distilled TabTransformer architectures. Subsequently, based on a curated domain-specific dataset encompassing 383 validated selenium nanoparticle records, systematic feature statistics, preprocessing pipelines, and predictive models were developed.

Empirical results demonstrate that on large-scale heterogeneous tabular data, TabTransformer achieved the highest $R^2$ (0.768) in regression tasks, whereas tree-based ensembles exhibited superior stability in classification tasks, with XGBoost yielding an area under the receiver operating characteristic curve (AUC) of 0.945. In the strict selenium nanoparticle reconstruction, 434 records from 24 publication groups were evaluated with five-fold outer and three-fold inner grouped nested cross-validation. The no-IC50 Random Forest baseline achieved an OOF MAE of 19.11\%, RMSE of 24.12\%, and $R^2=0.109$; XGBoost achieved MAE 19.49\%, RMSE 23.97\%, and $R^2=0.120$. Chemistry-enhanced features and an explicit log-dose--log-time interaction did not produce stable improvements. The IC50-assisted XGBoost configuration yielded the best numerical metrics (MAE 18.35\%, RMSE 23.25\%, $R^2=0.172$), but it is treated only as a response-correlated auxiliary reference because its paired improvement interval crossed zero.

These findings show that historical random-split performance should not be used as the prospective accuracy estimate for this small, heterogeneous archive. The reproducible data and evaluation workflow provides a basis for frozen-model external validation, but complete modifier categories, molecular weights, formulation identifiers, and independent wet-lab dose--response measurements remain necessary before deployment claims can be made.
